\documentclass[aspectratio=169,11pt]{beamer}

% ============================================
% PACKAGES
% ============================================
\usepackage[utf8]{inputenc}
\usepackage[T1]{fontenc}
\usepackage[ngerman]{babel}
\usepackage{lmodern}
\usepackage{tcolorbox}
\tcbuselibrary{skins, hooks}
\usepackage{listings}
\usepackage{xcolor}
\usepackage{fontawesome5}
\usepackage{tikz}
\usetikzlibrary{trees, shapes}

% Modern Font
\IfFileExists{FiraSans.sty}{\usepackage[sfdefault]{FiraSans}}{}
\renewcommand{\familydefault}{\sfdefault}

% ============================================
% COLORS & DESIGN SYSTEM
% ============================================
\definecolor{BgColor}{HTML}{FAFAFA}
\definecolor{Primary}{HTML}{2B3A42}
\definecolor{Accent}{HTML}{E84A5F}
\definecolor{Secondary}{HTML}{3F5765}
\definecolor{CodeBg}{HTML}{F0F0F0}
\definecolor{Success}{HTML}{28A745}

\setbeamercolor{background canvas}{bg=BgColor}
\setbeamercolor{title}{fg=Primary}
\setbeamercolor{frametitle}{bg=Primary, fg=white}
\setbeamercolor{structure}{fg=Accent}
\setbeamercolor{normal text}{fg=Primary}
\setbeamercolor{itemize item}{fg=Accent}

\setbeamertemplate{navigation symbols}{}
\setbeamertemplate{footline}{%
  \leavevmode%
  \hbox{%
  \begin{beamercolorbox}[wd=.333333\paperwidth,ht=2.25ex,dp=1ex,center]{author in head/foot}%
    \usebeamerfont{author in head/foot}\tiny Falko Himmel, Malte Baumann, Max Lehm
  \end{beamercolorbox}%
  \begin{beamercolorbox}[wd=.333333\paperwidth,ht=2.25ex,dp=1ex,center]{title in head/foot}%
    \usebeamerfont{title in head/foot}\tiny Grundlagen der Praktischen Informatik
  \end{beamercolorbox}%
  \begin{beamercolorbox}[wd=.333333\paperwidth,ht=2.25ex,dp=1ex,right]{date in head/foot}%
    \usebeamerfont{date in head/foot}\insertframenumber{} / \inserttotalframenumber\hspace*{2ex} 
  \end{beamercolorbox}}%
  \vskip0pt%
}

% ============================================
% CUSTOM BOXES
% ============================================
\tcbset{
    modernbox/.style={
        enhanced,
        boxrule=0pt,
        frame hidden,
        colback=white,
        coltitle=white,
        colbacktitle=Secondary,
        fonttitle=\bfseries,
        drop lifted shadow,
        arc=4pt,
        toptitle=2mm,
        bottomtitle=2mm,
        left=2mm,
        right=2mm
    },
    alertbox/.style={
        modernbox,
        colbacktitle=Accent
    }
}

% ============================================
% CODE LISTINGS
% ============================================
\lstdefinelanguage{Haskell}{
  keywords={class, data, deriving, do, else, if, in, import, let, module, of, then, type, where, case, error},
  keywordstyle=\color{Accent}\bfseries,
  identifierstyle=\color{Primary},
  sensitive=true,
  comment=[l]{--},
  commentstyle=\color{gray}\ttfamily\itshape,
  stringstyle=\color{teal}\ttfamily,
  morestring=[b]"
}

\lstset{
    language=Haskell,
    backgroundcolor=\color{CodeBg},
    basicstyle=\ttfamily\scriptsize,
    breaklines=true,
    frame=single,
    rulecolor=\color{CodeBg},
    frameround=tttt,
    numbers=left,
    numberstyle=\tiny\color{gray},
    xleftmargin=15pt,
    tabsize=4,
    showstringspaces=false
}

% ============================================
% TITLE PAGE
% ============================================
\title{\textbf{Tutorium 9}}
\subtitle{Altklausuraufgaben: Bäume}
\institute{Grundlagen der Praktischen Informatik}
\author{}
\date{\today}

\begin{document}

\begin{frame}[plain]
    \titlepage
\end{frame}

% ============================================
% WIEDERHOLUNG BÄUME
% ============================================
\section{Wiederholung}

\begin{frame}[fragile]{Binärbäume}
    Zur Erinnerung: Ein Baum und die \texttt{foldTree} Funktion.
    
    \begin{tcolorbox}[modernbox, title=Definition]
\begin{lstlisting}
data BinTree a = Node a (BinTree a) (BinTree a) | Empty
    deriving (Eq, Show)

foldTree :: (a -> b -> b -> b) -> b -> BinTree a -> b
foldTree fNode fEmpty = fold
    where
        fold (Node a l r) = fNode a (fold l) (fold r)
        fold Empty = fEmpty
\end{lstlisting}
    \end{tcolorbox}
\end{frame}

% ============================================
% ALTKLAUSUR
% ============================================

\begin{frame}[fragile]{Beispielbaum aus der Altklausur}
    \begin{tcolorbox}[modernbox, title=\texttt{bsp} Baum]
\begin{lstlisting}
bsp = Node 1 (Node 2 Empty (Node 3 Empty Empty)) 
             (Node 4 Empty Empty)
\end{lstlisting}
    \end{tcolorbox}
    
    \begin{center}
    \begin{tikzpicture}[level distance=1cm, sibling distance=2cm, every node/.style={circle, draw, minimum size=0.6cm}]
    \node {1}
      child {node {2}
        child {node[draw=none] {}}
        child {node {3}}
      }
      child {node {4}};
    \end{tikzpicture}
    \end{center}
\end{frame}

\section{Altklausuraufgaben (2024)}

\begin{frame}[fragile]{Aufgabe 4a: Innere Knoten zählen}
    \textbf{Aufgabe:} \texttt{countInnerNodes} zählt die Anzahl der inneren Knoten (Knoten, die kein Blatt sind).
    
    \begin{tcolorbox}[modernbox, title=Rekursive Lösung]
\begin{lstlisting}
countInnerNodes' :: BinTree a -> Int
countInnerNodes' Empty = 0 
countInnerNodes' (Node _ Empty Empty) = 0 
countInnerNodes' (Node _ l r) = 
    1 + countInnerNodes' l + countInnerNodes' r
\end{lstlisting}
    \end{tcolorbox}
\end{frame}

\begin{frame}[fragile]{Aufgabe 4b: Innere Knoten auflisten}
    \textbf{Aufgabe:} \texttt{listInnerNodes} listet die Werte aller inneren Knoten in einer Liste auf.
    
    \begin{tcolorbox}[modernbox, title=Rekursive Lösung]
\begin{lstlisting}
listInnerNodes' :: BinTree a -> [a]
listInnerNodes' Empty = []
listInnerNodes' (Node _ Empty Empty) = []
listInnerNodes' (Node a l r) = 
    a : (listInnerNodes' l ++ listInnerNodes' r)
\end{lstlisting}
    \end{tcolorbox}
\end{frame}

\begin{frame}[fragile]{Aufgabe 4c: Baum-Produkt mit \texttt{foldTree}}
    \textbf{Aufgabe:} \texttt{productTree} berechnet das Produkt aller Werte mit \texttt{foldTree}.
    
    \begin{tcolorbox}[modernbox, title=Vergleich: Rekursion vs. Fold]
\begin{lstlisting}
-- Rekursiv:
productTree' Empty = 1
productTree' (Node a l r) = a * productTree' l * productTree' r

-- Mit foldTree:
productTree'' :: Num a => BinTree a -> a
productTree'' = foldTree (\a l r -> a * l * r) 1
\end{lstlisting}
    \end{tcolorbox}
\end{frame}

% ============================================
% LIVE DEMO
% ============================================
\section{Praxis}

\begin{frame}[plain]
    \begin{center}
        \Huge \textbf{Live Demo} \\
        \vspace{0.5cm}
        \large \faLaptopCode~ \texttt{01\_altkausuraufgaben\_baeume.hs} \\
        \vspace{0.3cm}
        \normalsize Wir besprechen diese Klausuraufgaben zusammen in der IDE!
    \end{center}
\end{frame}

\end{document}
